%% Insert after the COMPLETE Section 2, not inside the fuel model.
\section{Proposed risk-aware two-stage trajectory planning framework}
\label{sec:method}

The framework separates route selection, geometric refinement and speed
assignment while retaining a common mission evaluator. The coarse search
first identifies a route under the normalized time--exposure--fuel objective.
A fine search then resolves its local geometry within a three-dimensional
corridor. Cubic B-splines replace the resulting polyline with differentiable
candidates, which are screened against geometric constraints and ranked
using a dimensionless refinement score. Finally, a forward--backward pass
assigns speeds along the selected curve, and the mission metrics are
recomputed under the resulting motion. This decomposition limits the size
of the search state while making the distinction between search estimates
and final trajectory performance explicit.

%% FIGURE PLACEHOLDER: overall framework; use vector PDF for the final artwork.
%% Show mission/environment -> three raw anchors and shared scales -> coarse A*
%% -> 3-D corridor -> fine A* -> spline candidates/repairs -> geometric checks
%% -> minimum geometric score -> strict speed assignment -> final T, C, F, J.
%% Show rejection exits for search, geometry and speed failures. Add a dashed
%% coarse-to-spline branch labelled "72-case sensitivity protocol".
%% \begin{figure}[t]
%% \centering
%% \includegraphics[width=\linewidth]{figures/framework.pdf}
%% \caption{Planning and evaluation workflow. Mission preferences govern graph
%% search, whereas geometric refinement and speed assignment are subsequent
%% stages with explicit feasibility checks.}
%% \label{fig:framework}
%% \end{figure}

\subsection{Multi-objective risk-aware A* search}
\label{sec:astar}

The A* search~\cite{hart1968} represents position on an anisotropic Cartesian lattice
with horizontal and vertical spacings $\delta_{xy}$ and $\delta_z$.
Candidate displacements are drawn from
\begin{equation}
 \mathcal{D}=\{(a\delta_{xy},b\delta_{xy},c\delta_z):
 a,b,c\in\{-1,0,1\},\ (a,b,c)\ne(0,0,0)\}.
 \label{me:eq:neighbors}
\end{equation}
Purely vertical moves are excluded, and a displacement is retained only
if its climb angle satisfies the prescribed bound. The map boundary,
altitude envelope and hard exclusion regions further restrict the
admissible graph. Each retained edge is checked at ten spatial locations,
including its endpoints; the nine intervening subsegments supply the
midpoint quadrature for exposure and fuel. A graph state contains position
only. Heading, curvature and acceleration feasibility are consequently
resolved after search rather than represented as additional state variables.

For a partial path ending at $\mathbf p$, the accumulated search cost is
$G(\mathbf p)=\sum_e c_e$, with $c_e$ defined in
Eq.~\eqref{pr:eq:edgeobjective}. The A* priority is
\begin{equation}
 f(\mathbf p)=G(\mathbf p)+h(\mathbf p),\qquad
 h(\mathbf p)=\frac{w_T}{\Delta_T}
              \frac{\|\mathbf p-\mathbf p_g\|}{v_{\mathrm{cr}}}.
 \label{me:eq:heuristic}
\end{equation}
No heuristic inflation is used. The heuristic retains only a lower bound
on the time contribution: risk and fuel are nonnegative, and the search
time of an edge is at least its length divided by $v_{\mathrm{cr}}$.
For adjacent states $\mathbf p$ and $\mathbf q$, the triangle inequality
therefore gives
\begin{equation}
 h(\mathbf p)-h(\mathbf q)
 \le \frac{w_T\|\mathbf p-\mathbf q\|}{\Delta_Tv_{\mathrm{cr}}}
 \le c_{(\mathbf p,\mathbf q)}.
 \label{me:eq:consistency}
\end{equation}
Thus, the heuristic is consistent for the fixed graph and edge model.
When $w_T=0$, it becomes zero and the primary search reduces to uniform-cost
expansion. Equal primary priorities are ordered using accumulated geometric
length plus Euclidean distance to the goal, providing a deterministic
geometric preference in flat-cost regions. This secondary criterion does
not add a distance penalty to the mission objective.

The constant lower-reference offset in Eq.~\eqref{pr:eq:offset} is restored
when reporting the complete path objective. It is not included repeatedly
in edge costs. The resulting minimum-cost interpretation concerns the
discretized search problem, subject to numerical tolerances and successful
search termination. It does not extend to the subsequently modified and
retimed continuous trajectory.
%% Cite a verified primary A* reference here when finalizing the bibliography.

\subsection{Coarse-to-fine three-dimensional planning}
\label{sec:coarsefine}

A uniform fine lattice over the entire workspace requires many expansions
in regions that are irrelevant to a particular mission preference. The
two-stage search instead uses a coarse path $\Gamma_c$ to restrict the
domain of a second, finer search. For any polyline $\Gamma$, define
\begin{equation}
 d(\mathbf p,\Gamma)=
 \min_{[\mathbf a,\mathbf b]\in\Gamma}\;
 \min_{\tau\in[0,1]}
 \|\mathbf p-\mathbf a-\tau(\mathbf b-\mathbf a)\|.
 \label{me:eq:distance}
\end{equation}
This distance uses the closest point on a segment, rather than the closest
waypoint, and treats horizontal and vertical displacement in the same
three-dimensional metric. The search corridor is
\begin{equation}
 \mathcal C(\Gamma_c,r_c)=
 \{\mathbf p\in\mathcal W:d(\mathbf p,\Gamma_c)\le r_c\}.
 \label{me:eq:corridor}
\end{equation}
Where grid snapping changes an endpoint, the physical endpoint is connected
to the coarse polyline when constructing the corridor reference. The
endpoints used in the present experiments lie on both lattices.

The fine graph retains nodes inside this corridor and samples candidate
edges to prevent a connection between admissible endpoints from leaving
the corridor unchecked. For the reported settings, corridor probes are
spaced no more than approximately $100~\mathrm{m}$ along an edge, with at
least three probes. Altitude and obstacle screening remain active. Both
search stages use the same objective weights and the same anchor-derived
normalization scales; the fine stage does not recompute its own anchors.

The fine result $\Gamma_f$ is obtained by a new graph search, not by
interpolating $\Gamma_c$. Its finer directional and altitude resolution
permits local route changes, but the corridor limits which alternatives
can be explored. The construction therefore trades search coverage for
computational tractability; it does not guarantee equivalence to an
unrestricted fine-grid search. A failed fine search is reported as a
failure rather than relabeling the coarse path as a fine result.

%% FIGURE PLACEHOLDER: coarse/corridor/fine geometry, preferably two panels.
%% Use the SAME actual case: (a) plan view with coarse and fine paths and
%% corridor; (b) 3-D detail or altitude profile showing vertical refinement.
%% Do not draw the 3-D corridor merely as a fixed-width horizontal strip.
%% \begin{figure}[t]
%% \centering
%% \includegraphics[width=\linewidth]{figures/coarse_corridor_fine.pdf}
%% \caption{Coarse-to-fine route refinement within a three-dimensional
%% distance corridor. The corridor restricts the fine search domain.}
%% \label{fig:coarsefine}
%% \end{figure}

\subsection{Cubic B-spline trajectory refinement}
\label{sec:bspline}

The fine polyline satisfies the edge-level search constraints but retains
direction discontinuities at its vertices. Let
$\{\mathbf P_i\}_{i=0}^{n-1}$ denote its ordered vertices after removal of
consecutive duplicates. A cubic parametric B-spline is fitted as
\begin{equation}
 \mathbf q_\sigma(u)=\sum_{j=0}^{K_\sigma-1}
 B_{j,3}(u)\mathbf c_j,\qquad u\in[0,1],
 \label{me:eq:spline}
\end{equation}
where $B_{j,3}$ are cubic basis functions, $\mathbf c_j$ are control points,
and $\sigma$ controls the fit. Chord-length parameterization assigns
$u_i=\sum_{k=0}^{i-1}\|\mathbf P_{k+1}-\mathbf P_k\|/L_{\mathrm{raw}}$,
where $L_{\mathrm{raw}}$ is the polyline length. The initial FITPACK
smoothing fit~\cite{dierckx1993} uses the weighted residual condition
\begin{equation}
 \sum_{i=0}^{n-1}\omega_i^2
       \|\mathbf P_i-\mathbf q_\sigma(u_i)\|^2\le\sigma,
 \qquad
 \omega_i=
 \begin{cases}
 10^5,&i\in\{0,n-1\},\\
 1,&\text{otherwise}.
 \end{cases}
 \label{me:eq:fit}
\end{equation}
Large endpoint weights promote endpoint retention, which is subsequently
verified with an explicit tolerance. The smoothing condition controls the
initial fit; it is not the final mission objective and is not assumed to
remain an exact residual bound after geometric repair.

To explore the balance between polyline fidelity and curvature reduction,
the smoothing budgets are scaled to the input resolution:
\begin{equation}
 \sigma_j=n(\delta_{xy}\beta_j)^2.
 \label{me:eq:budgets}
\end{equation}
The fixed set of 21 dimensionless factors is
\begin{equation}
 \begin{split}
 \mathcal B=\{&0.02,0.03,0.04,0.05,0.0625,0.075,0.09,\\
              &0.105,0.125,0.15,0.18,0.22,0.25,0.35,\\
              &0.50,0.75,1.0,1.5,2.0,3.0,4.0\}.
 \end{split}
 \label{me:eq:budgetset}
\end{equation}
Every budget is examined before selecting a candidate. Small budgets
retain the lattice geometry more closely; larger budgets allow greater
deformation but must still satisfy the local-deviation and flight-envelope
tests. The input is $\Gamma_f$ for the complete two-stage method and
$\Gamma_c$ for the separate coarse-search sensitivity protocol.

Geometric repair addresses constraint violations introduced by smoothing.
Altitude repair raises deficient sampled altitudes relative to the spatial
altitude floor and refits the corrected points. Climb repair limits the
vertical slope using margins below $\tan\gamma_{\max}$ while retaining
the sampled horizontal coordinates before refitting. These operations
generate revised candidates, each of
which is rechecked against the full geometric test set. Repair is a
feasibility procedure and need not reduce the mission objective.

The terminal heading is inferred from the incoming raw route over a
lookback length of $\max(5\delta_{xy},1000~\mathrm{m})$, using the horizontal
chord from the selected upstream vertex to the goal. If the initial
candidate pass finds no geometrically feasible curve and includes a
heading violation, a second pass aligns the final control-polygon edge
with this reference and reevaluates all smoothing budgets. The adjustment
retains the endpoint coordinates and control-point altitudes. The method
thus preserves the approach direction implied by the route; it does not
enforce an independently specified terminal mission heading.

For regular curves with simple interior knots, cubic B-splines provide
$C^2$ geometric continuity and permit curvature evaluation from analytic
derivatives. This property removes polyline corners, but does not by itself
establish feasibility or imply continuous jerk in the time-parameterized
motion. Those issues require the subsequent geometric and speed checks.
%% FITPACK interface details are recorded in literature_alignment.md.

\subsection{Geometric feasibility validation and candidate selection}
\label{sec:trajectory_validation}

Candidate curves are evaluated at $M=4096$ uniformly spaced parameter
values. Their analytic derivatives give spatial curvature through
Eq.~\eqref{pr:eq:curvature} and the local flight-path angle
\begin{equation}
 \gamma(u)=\arctan\!left(
 \frac{|q'_z(u)|}{\sqrt{q'_x(u)^2+q'_y(u)^2}}\right).
 \label{me:eq:splineangle}
\end{equation}
Finite coordinates and nonzero tangents are required. The geometric
turning-radius threshold is
\begin{equation}
 R_{\mathrm{ref}}=
 \max\left\{R_{\mathrm{floor}},
       \frac{v_{\mathrm{stall}}^2}{g\tan\phi_{\max}}\right\}.
 \label{me:eq:radius}
\end{equation}
Requiring $\kappa\le1/R_{\mathrm{ref}}$ prevents acceptance of a turn
that already exceeds the bank envelope at stall speed. It does not require
the curve to support cruise speed everywhere, because speed can subsequently
decrease in curved sections.

A candidate must also satisfy the altitude envelope, avoid the hard
exclusion regions, retain both endpoints within $2~\mathrm{m}$, and match
the terminal horizontal heading within $12^\circ$. Its maximum distance
from the input polyline must satisfy
\begin{equation}
 d(\mathbf q_\sigma(u),\Gamma_{\mathrm{raw}})\le r_s
 \label{me:eq:locality}
\end{equation}
at the validation samples. Here $r_s$ is a refinement locality limit,
distinct from the coarse-search corridor radius $r_c$. In the default
configuration, membership of the final spline in the coarse corridor is
recorded diagnostically, whereas locality relative to the raw input curve
is enforced. Obstacles and the altitude envelope remain hard requirements.

Among the geometrically admissible candidates, selection balances bending,
curvature variation, deviation and length. Let $s$ denote arc length,
$L_\sigma$ the candidate length, and $\kappa_s=\mathrm d\kappa/\mathrm ds$.
The dimensionless components are
\begin{align}
 Q_b&=\frac{1}{L_\sigma}
       \int_0^{L_\sigma}(R_{\mathrm{ref}}\kappa)^2\,\mathrm ds,
       \label{me:eq:bending}\\
 Q_v&=R_{\mathrm{ref}}^3
       \int_0^{L_\sigma}\kappa_s^2\,\mathrm ds,
       \label{me:eq:variation}\\
 Q_d&=\frac{1}{r_s}
       \sqrt{\frac{1}{M}\sum_{m=0}^{M-1}
       d(\mathbf q_\sigma(u_m),\Gamma_{\mathrm{raw}})^2},
       \qquad u_m=\frac{m}{M-1},
       \label{me:eq:deviation}\\
 Q_l&=\frac{L_\sigma}{L_{\mathrm{raw}}}.
       \label{me:eq:length}
\end{align}
The integrals use the cumulative chord-length coordinate of the dense
samples and trapezoidal quadrature; $\kappa_s$ is estimated numerically
on that coordinate. In contrast, $Q_d$ is the root-mean-square deviation
over uniform parameter samples, not an arc-length-weighted integral.
The curvature-variation term penalizes spatial changes in curvature;
it is not a roll-rate constraint.

The refinement score and selection rule are
\begin{equation}
 J_s=0.4Q_b+0.3Q_v+0.2Q_d+0.1Q_l,\qquad
 \mathbf q^*=\mathop{\arg\min}_{\mathbf q\in\mathcal Q_{\mathrm{geom}}}
 J_s(\mathbf q),
 \label{me:eq:selection}
\end{equation}
where $\mathcal Q_{\mathrm{geom}}$ contains the admitted candidates from
the applicable fitting and repair pass. The four coefficients are fixed
geometric design settings, independent of $(w_T,w_C,w_F)$. Selection
enumerates the candidate set rather than stopping at the first feasible
smoothing budget. It is a finite candidate comparison, not a global
optimization over all spline control points.

If no geometrically admissible candidate exists, refinement fails. A
diagnostic fallback curve is not accepted as a mission trajectory. The
selected curve must next pass strict speed validation; the present
procedure does not rerank alternative candidates following a speed-stage
failure. Final acceptance therefore requires both stages to succeed.
Because geometric tests use finite sampling, acceptance has the numerical
meaning stated in Section~\ref{pr:sec:constraints}.

\subsection{Adaptive speed planning and final mission evaluation}
\label{sec:speed}

The selected spline is resampled at approximately uniform arc-length
spacing. A dense set of 8000 parameter probes supplies the cumulative
length map, which is inverted to obtain the final samples
$\{\mathbf q_i\}_{i=0}^{N}$. Their chord lengths are
$\ell_i=\|\mathbf q_{i+1}-\mathbf q_i\|$, and $\kappa_i$ is evaluated
from the spline derivatives at the corresponding parameters. The target
spacing is $100~\mathrm{m}$ in both experimental protocols.

The vertical-rate ceiling of a nondegenerate segment is
\begin{equation}
 \bar v_{z,i}=\begin{cases}
 v_{\mathrm{asc},\max}\ell_i/|z_{i+1}-z_i|,&z_{i+1}>z_i,\\
 v_{\mathrm{des},\max}\ell_i/|z_{i+1}-z_i|,&z_{i+1}<z_i,\\
 +\infty,&z_{i+1}=z_i.
 \end{cases}
 \label{me:eq:verticalcap}
\end{equation}
This ceiling is assigned to both endpoints of the segment. At node $i$,
the admissible speed ceiling is
\begin{equation}
 \bar v_i=\min\left\{
 v_{\mathrm{cr}},v_{\max},
 \sqrt{\frac{g\tan\phi_{\max}}{\kappa_i}},
 \min_{e\in\mathcal E_i}\bar v_{z,e}\right\},
 \label{me:eq:nodecap}
\end{equation}
where $\mathcal E_i$ comprises adjacent segments and the curvature term is
infinite on a straight segment. A ceiling below $v_{\mathrm{stall}}$
causes rejection, rather than a speed assignment below stall.

With $v_0^+=v_{\mathrm{stall}}$, the forward pass propagates the acceleration
bound:
\begin{equation}
 v_{i+1}^+=\min\left\{\bar v_{i+1},
       \sqrt{(v_i^+)^2+2a_{\mathrm{acc},\max}\ell_i}\right\},
 \qquad i=0,\ldots,N-1.
 \label{me:eq:forward}
\end{equation}
The backward pass anticipates downstream speed reductions:
\begin{equation}
 v_N=v_N^+,\qquad
 v_i=\min\left\{v_i^+,
       \sqrt{v_{i+1}^2+2a_{\mathrm{dec},\max}\ell_i}\right\},
 \qquad i=N-1,\ldots,0.
 \label{me:eq:backward}
\end{equation}
No terminal-speed equality is introduced. Applying a segment's vertical
ceiling at both ends bounds its largest vertical speed under the adopted
constant-tangential-acceleration segment model. The final validation checks
speed, turning, climb angle, vertical rate and tangential acceleration;
failure prevents the trajectory from being reported as feasible.

The accepted speed sequence supplies segment times through
Eq.~\eqref{pr:eq:segmenttime}. Time, exposure and fuel are then evaluated
with the common subsegment quadrature in Eq.~\eqref{pr:eq:quadrature},
including renewed altitude and collision screening, and normalized using
the unchanged anchor scales. This last evaluation is essential because
spline deformation changes the spatial threat field sampled by the
vehicle, while retiming changes both residence time and fuel expenditure.
Consequently, a lower geometric refinement score does not by itself imply
a lower final mission objective. The reported output is the accepted
geometry, its assigned speed profile and the recomputed $(T,C,F,J)$;
search-stage estimates are retained separately for diagnosis and stage comparisons.

%% FIGURE PLACEHOLDER: use measured outputs in Section 4.2, not synthetic data.
%% Suggested panels: local raw/spline geometry; analytic spline curvature
%% with 1/R_ref; speed with stall/local ceilings; climb and vertical-rate
%% profiles with limits. Raw-polyline corner curvature is diagnostic only.
